\documentclass[12pt]{article}
\usepackage{amsmath,amssymb,amsfonts,amsthm}
\usepackage[margin=1in]{geometry}

\begin{document}

\begin{center}
{\Large\bfseries Chapter 4, Problem 17}\\[6pt]
Byron \& Fuller, \textit{Mathematics of Classical and Quantum Physics}
\end{center}

\bigskip

\textbf{Problem.} The projection operator $P_i$ is defined as follows:
\[
P_i x = x_i(x_i, x),
\]
where $x_i$ is a unit vector. This operator is called a projection operator because all of $x$ are in the direction of $x_i$ and the length of $x'$ equals the component of $x$ in $x_i$ direction, namely $(x_i, x)$. In the following, let the $x_i$ be a set of orthonormal vectors which span the $n$-dimensional space $V$. Prove:
\begin{itemize}
\item[(a)] $P_i$ is idempotent.
\item[(b)] $P_i P_j = 0$ for $i \neq j$. Interpret geometrically.
\item[(c)] $P_i$ has no inverse.
\item[(d)] $\sum_i P_i = I$ (the identity operator).
\item[(e)] Let $A$ be a self-adjoint linear operator defined on an $n$-dimensional vector space $V$ with $n$ eigenvalues $e_i$ and $n$ eigenvectors $x_i$. Prove that $A$ may be written as
\[
A = \sum_i e_i P_i,
\]
where $P_i$ is defined above. This is known as the spectral theorem for self-adjoint linear operators. (The analogous theorem for normal operators can also be easily proved.) Note the consistency of this result with parts (d) and (e) of this problem.
\end{itemize}

\bigskip
\textbf{Solution.}

\textbf{(a) $P_i$ is idempotent ($P_i^2 = P_i$):}

For any vector $x$:
\[
P_i^2 x = P_i(P_i x) = P_i\bigl(x_i(x_i, x)\bigr) = (x_i, x)\,P_i(x_i) = (x_i, x)\,x_i(x_i, x_i).
\]
Since $x_i$ is a unit vector, $(x_i, x_i) = 1$, so:
\[
P_i^2 x = (x_i, x)\,x_i = P_i x.
\]
Therefore $P_i^2 = P_i$. \qed

\bigskip
\textbf{(b) $P_i P_j = 0$ for $i \neq j$:}

\[
P_i P_j x = P_i\bigl(x_j(x_j, x)\bigr) = (x_j, x)\,P_i(x_j) = (x_j, x)\,x_i(x_i, x_j).
\]
Since $\{x_i\}$ are orthonormal, $(x_i, x_j) = 0$ for $i \neq j$. Therefore $P_i P_j = 0$.

\textbf{Geometric interpretation:} Projecting first onto the $x_j$ direction and then onto the $x_i$ direction (with $i \neq j$) gives zero because these directions are orthogonal. A vector in the $x_j$ direction has no component along $x_i$.

\bigskip
\textbf{(c) $P_i$ has no inverse:}

Since $P_i^2 = P_i$, if $P_i$ had an inverse $P_i^{-1}$, we could write:
\[
P_i^{-1}P_i^2 = P_i^{-1}P_i \implies P_i = I,
\]
but $P_i \neq I$ for $n > 1$ (it projects onto a one-dimensional subspace). Contradiction.

Alternatively, $P_i$ maps the entire orthogonal complement of $x_i$ to zero, so its kernel is nontrivial (dimension $n-1$), and hence $P_i$ is singular.

\bigskip
\textbf{(d) $\sum_i P_i = I$:}

For any vector $x$, since $\{x_i\}$ is a complete orthonormal basis:
\[
x = \sum_{i=1}^n (x_i, x)\,x_i = \sum_{i=1}^n P_i x.
\]
Therefore $\left(\sum_i P_i\right)x = x$ for all $x$, which means $\sum_i P_i = I$. \qed

\bigskip
\textbf{(e) Spectral decomposition $A = \sum_i e_i P_i$:}

Since $A$ is self-adjoint on an $n$-dimensional space, it has $n$ eigenvalues $e_i$ (real) with corresponding orthonormal eigenvectors $x_i$: $Ax_i = e_i x_i$.

For any vector $v$, expand in the eigenbasis:
\[
v = \sum_j (x_j, v)\,x_j.
\]
Applying $A$:
\[
Av = \sum_j (x_j, v)\,Ax_j = \sum_j e_j(x_j, v)\,x_j = \sum_j e_j P_j v.
\]
Since this holds for all $v$:
\[
A = \sum_i e_i P_i. \qed
\]

\textbf{Consistency check:} From (d), $I = \sum_i P_i$. Acting with $A$: $A = AI = A\sum_i P_i = \sum_i AP_i$. Since $AP_i v = A\bigl(x_i(x_i,v)\bigr) = (x_i,v)Ax_i = e_i(x_i,v)x_i = e_i P_i v$, we get $A = \sum_i e_i P_i$, consistent. \qed

\end{document}
